\documentclass[12pt,a4paper]{article}

% ============================================================
% PACKAGES
% ============================================================

\usepackage[a4paper,margin=1in]{geometry}
\usepackage{graphicx}
\usepackage{float}
\usepackage{placeins}
\usepackage{tabularx}
\usepackage{array}
\usepackage{enumitem}
\usepackage{xcolor}
\usepackage{hyperref}
\usepackage{setspace}
\usepackage{fancyhdr}
\usepackage{titlesec}

% ============================================================
% COLORS
% ============================================================

\definecolor{linkblue}{RGB}{0,90,170}

% ============================================================
% HYPERLINK SETTINGS
% ============================================================

\hypersetup{
    colorlinks=true,
    linkcolor=linkblue,
    urlcolor=linkblue,
    citecolor=linkblue
}

% ============================================================
% DOCUMENT SETTINGS
% ============================================================

\setstretch{1.08}

\setlength{\parindent}{0pt}
\setlength{\parskip}{6pt}

\setlist[itemize]{
    leftmargin=1.5em
}

\setlist[enumerate]{
    leftmargin=1.7em
}

% ============================================================
% HEADING FORMATTING
% ============================================================

\titleformat{\section}
    {\Large\bfseries}
    {\thesection}
    {0.7em}
    {}

\titleformat{\subsection}
    {\large\bfseries}
    {\thesubsection}
    {0.7em}
    {}

\titleformat{\subsubsection}
    {\normalsize\bfseries}
    {\thesubsubsection}
    {0.7em}
    {}

% ============================================================
% HEADER AND FOOTER
% ============================================================

\pagestyle{fancy}
\fancyhf{}

\fancyhead[L]{%
    \small\textbf{InternNova -- Week 3}
}

\fancyhead[R]{%
    \small\textbf{Malware, Phishing \& Security Awareness}
}

\fancyfoot[C]{%
    \thepage
}

\renewcommand{\headrulewidth}{0.4pt}
\renewcommand{\footrulewidth}{0pt}

% ============================================================
% REFERENCE COMMAND
% ============================================================

\newcommand{\refs}[3]{%
    \textbf{Reference Basis:} Authoritative cybersecurity guidance\par
    \textbf{Reference Paper / Standard:} Relevant guidance or standard\par
    \textbf{Learning Video:} Educational cybersecurity video\par
    \href{#1}{\textcolor{linkblue}{View Reference}}
    \quad
    \href{#2}{\textcolor{linkblue}{View Standard}}
    \quad
    \href{#3}{\textcolor{linkblue}{Watch on YouTube}}
}

% ============================================================
% DOCUMENT
% ============================================================

\begin{document}

% ============================================================
% COVER PAGE
% ============================================================


% ---------- Cover ----------
\begin{titlepage}
\centering
\vspace*{1.5cm}

{\Large\textbf{InternNova Internship Program}\par}
\vspace{0.8cm}

{\large\textbf{Week 3 Assignment}\par}
\vspace{1.2cm}

{\LARGE\textbf{Malware, Phishing \& Security Awareness}\par}
\vspace{1.0cm}

{\Large\textbf{Task 1 \& Task 2}\par}
\vspace{0.6cm}

{\large\textbf{Malware, Phishing \& Security Awareness}\par}

\vfill

\begin{tabular}{>{\bfseries}l l}
Intern Name: & NAJMUDEEN S\\
Internship: & InternNova\\
Week: & 3\\
Assignment: & Malware, Phishing \& Security Awareness\\
Month: & October 2026\\
\end{tabular}

\vfill

\end{titlepage}

\clearpage

% ============================================================
% TASK 1
% ============================================================

\section{Task 1: Understanding Malware}

\subsection{Introduction}

Malware is short for malicious software. It refers to software designed to perform unauthorized or harmful actions on a computer, mobile device, network, or information system. Different malware types behave differently: some infect files, some spread between systems, some disguise themselves as legitimate software, some restrict access to data, and some secretly collect information.

% ============================================================
% VIRUS
% ============================================================

\subsection{Virus}

\subsubsection{Definition}

A virus is malware that attaches itself to a file or program and can replicate when the infected content is executed.

\subsubsection{How It Works}

A typical virus becomes active when an infected file is executed. The malicious code can then modify or attach itself to other files, allowing the infection to continue when those files are opened or shared.

\begin{figure}[H]
\centering
\includegraphics[width=.58\textwidth]{images/virus.png}
\caption{Virus Infection and Replication}
\end{figure}

\subsubsection{Real-World Example}

CIH, also known as the Chernobyl virus, is a historical example of a computer virus associated with infected executable files and destructive activity.

\subsubsection{Cybersecurity Connection}

Viruses are relevant to endpoint security because they can modify legitimate files and affect system integrity. Analysts may investigate antivirus detections, unusual file modifications, suspicious processes, and unexpected system behavior.

\subsubsection{Security Risk and Protection}

\textbf{Potential Risk:} An infected file can spread malicious code when it is executed or shared.

\textbf{Protection:}

\begin{itemize}
    \item Use reputable security software.
    \item Keep systems updated.
    \item Avoid unexpected attachments.
    \item Download software from trusted sources.
    \item Scan suspicious files.
\end{itemize}

\refs
{https://www.cisa.gov/news-events/news/malware}
{https://www.cisa.gov/news-events/news/malware}
{https://www.youtube.com/results?search_query=Professor+Messer+malware+virus}

% ============================================================
% WORM
% ============================================================

\subsection{Worm}

\subsubsection{Definition}

A worm is malware that can replicate itself and spread from one system to another without requiring the user to manually copy it.

\subsubsection{How It Works}

A worm may search for reachable or vulnerable systems and attempt to spread automatically. An infected system can then continue the propagation process.

\begin{figure}[H]
\centering
\includegraphics[width=.58\textwidth]{images/worm.png}
\caption{Worm Self-Propagation}
\end{figure}

\subsubsection{Real-World Example}

Conficker is a well-known computer worm that spread through Windows systems by exploiting vulnerabilities and weak configurations.

\subsubsection{Cybersecurity Connection}

Worms are important in network security because they can spread rapidly. Analysts may monitor unusual network traffic, repeated connection attempts, vulnerability exploitation, and unexpected endpoint behavior.

\subsubsection{Security Risk and Protection}

\textbf{Potential Risk:} A worm can automatically spread across vulnerable systems.

\textbf{Protection:}

\begin{itemize}
    \item Apply security updates.
    \item Use firewalls.
    \item Disable unnecessary services.
    \item Segment networks where appropriate.
    \item Monitor unusual network activity.
\end{itemize}

\refs
{https://www.cisa.gov/news-events/news/malware}
{https://www.cisa.gov/news-events/news/malware}
{https://www.youtube.com/results?search_query=Professor+Messer+computer+worm+malware}

% ============================================================
% TROJAN
% ============================================================

\subsection{Trojan}

\subsubsection{Definition}

A Trojan is malware that disguises itself as legitimate or useful software to trick a user into installing or executing it.

\subsubsection{How It Works}

A Trojan may be distributed through a fake application, document, attachment, or apparently legitimate file. After execution, it can perform unauthorized actions.

\begin{figure}[H]
\centering
\includegraphics[width=.58\textwidth]{images/trojan.png}
\caption{Trojan Deceptive Software Process}
\end{figure}

\subsubsection{Real-World Example}

Zeus is a well-known example of malware associated with Trojan-based distribution and theft of sensitive information such as banking credentials.

\subsubsection{Cybersecurity Connection}

Trojans are important because they often depend on deception or social engineering. Analysts may investigate unexpected installations, suspicious processes, unusual connections, and endpoint alerts.

\subsubsection{Security Risk and Protection}

\textbf{Potential Risk:} A user may unknowingly install malicious software because it appears legitimate.

\textbf{Protection:}

\begin{itemize}
    \item Use trusted software sources.
    \item Avoid pirated or modified software.
    \item Verify software publishers.
    \item Use endpoint security.
    \item Avoid unexpected attachments.
\end{itemize}

\refs
{https://www.cisa.gov/news-events/news/malware}
{https://www.cisa.gov/news-events/news/malware}
{https://www.youtube.com/results?search_query=Professor+Messer+Trojan+malware}

% ============================================================
% RANSOMWARE
% ============================================================

\subsection{Ransomware}

\subsubsection{Definition}

Ransomware is malware that can prevent access to files or systems, commonly by encrypting data, and then demands payment from the victim.

\subsubsection{How It Works}

Ransomware may enter through a malicious attachment, compromised account, vulnerable service, or another initial access method. It may identify files, encrypt them, and display a ransom demand.

\begin{figure}[H]
\centering
\includegraphics[width=.58\textwidth]{images/ransomware.png}
\caption{Ransomware File Encryption and Ransom Demand}
\end{figure}

\subsubsection{Real-World Example}

WannaCry is a well-known ransomware outbreak that affected computer systems in multiple countries.

\subsubsection{Cybersecurity Connection}

Ransomware can affect data availability and operations. Security teams may monitor endpoint alerts, file activity, authentication events, backups, and network traffic.

\subsubsection{Security Risk and Protection}

\textbf{Potential Risk:} Important files may become inaccessible and operations may be disrupted.

\textbf{Protection:}

\begin{itemize}
    \item Maintain protected backups.
    \item Patch systems.
    \item Use endpoint security.
    \item Restrict unnecessary network access.
    \item Train users to recognize suspicious links and attachments.
\end{itemize}

\refs
{https://www.cisa.gov/stopransomware}
{https://www.cisa.gov/stopransomware/ransomware-guide}
{https://www.youtube.com/results?search_query=Professor+Messer+ransomware}

% ============================================================
% SPYWARE
% ============================================================

\subsection{Spyware}

\subsubsection{Definition}

Spyware is malware designed to operate without the user's knowledge and collect information about user activity or sensitive information.

\subsubsection{How It Works}

Spyware may run in the background after installation. Depending on its capabilities, it can collect browsing activity, credentials, or other sensitive data and transmit it to an unauthorized party.

\begin{figure}[H]
\centering
\includegraphics[width=.58\textwidth]{images/spyware.png}
\caption{Spyware Information Collection}
\end{figure}

\subsubsection{Real-World Example}

Pegasus is an example of sophisticated spyware associated with targeted surveillance of mobile devices.

\subsubsection{Cybersecurity Connection}

Spyware is relevant to endpoint security and privacy because unauthorized collection can expose sensitive data. Analysts may investigate unusual processes, applications, connections, and endpoint alerts.

\subsubsection{Security Risk and Protection}

\textbf{Potential Risk:} Sensitive information may be collected without the user's knowledge.

\textbf{Protection:}

\begin{itemize}
    \item Install trusted applications.
    \item Keep systems updated.
    \item Use reputable security software.
    \item Review applications and permissions.
    \item Avoid suspicious links and downloads.
\end{itemize}

\refs
{https://www.cisa.gov/news-events/news/malware}
{https://www.cisa.gov/news-events/news/malware}
{https://www.youtube.com/results?search_query=Professor+Messer+spyware+malware}

% ============================================================
% OVERALL MALWARE DIAGRAMS
% ============================================================


\subsection{Overall Malware Diagrams}

\begin{figure}[H]
\centering
\includegraphics[width=.68\textwidth]{images/malware_lifecycle.png}
\caption{General Malware Infection Lifecycle}
\end{figure}

\begin{figure}[H]
\centering
\includegraphics[width=.95\textwidth]{images/malware_types_overview.png}
\caption{Common Malware Types Overview}
\end{figure}

% ============================================================
% MALWARE COMPARISON
% ============================================================

\subsection{Malware Comparison}

\begin{table}[H]
\centering
\small
\renewcommand{\arraystretch}{1.25}

\begin{tabularx}{\textwidth}{
|p{2.2cm}|X|X|X|
}
\hline
\textbf{Malware}
&
\textbf{Main Behavior}
&
\textbf{Typical Impact}
&
\textbf{Basic Protection}
\\
\hline

Virus
&
Infects files and replicates when activated
&
File/system modification
&
Updates and security software
\\
\hline

Worm
&
Replicates and spreads between systems
&
Rapid network/system infection
&
Patching, firewalls, segmentation
\\
\hline

Trojan
&
Disguises malicious software as legitimate software
&
Unauthorized activity
&
Trusted software sources
\\
\hline

Ransomware
&
Encrypts or restricts access to data
&
Loss of data availability
&
Backups and patching
\\
\hline

Spyware
&
Secretly collects information
&
Privacy/information loss
&
Trusted apps and endpoint security
\\
\hline

\end{tabularx}

\caption{Comparison of Common Malware Types}

\end{table}

% ============================================================
% MALWARE MINI PRACTICAL
% ============================================================

\subsection{Mini Practical / Try It Yourself}

A safe malware-awareness activity can be performed without downloading or executing real malware.

\begin{enumerate}
    \item Open the security settings of your operating system.
    \item Check whether real-time protection is enabled.
    \item Check the latest security update.
    \item Review installed applications.
    \item Verify unfamiliar applications before taking action.
\end{enumerate}

\textbf{Safety Note:} Do not download, execute, or interact with real malware for this assignment.

% ============================================================
% MALWARE CYBERSECURITY CONNECTION
% ============================================================

\subsection{Cybersecurity Connection}

Malware knowledge is important for cybersecurity analysts because malware can affect endpoints, networks, applications, and data. Analysts may use endpoint detection tools, security alerts, logs, network traffic, and file activity to investigate suspicious behavior.

% ============================================================
% MALWARE SECURITY RISKS
% ============================================================

\subsection{Security Risks and Protection}

\textbf{Risks:} Unauthorized access, information theft, file destruction, loss of data availability, and spread across connected systems.

\textbf{Protection:} Keep systems updated; use reputable security software; maintain protected backups; use strong authentication; use trusted software sources; avoid suspicious attachments and links; monitor unusual activity.

% ============================================================
% MALWARE QUICK CHECK
% ============================================================

\subsection{Quick Check}

\begin{enumerate}
    \item Which malware attaches to files?
    \item Which malware spreads automatically?
    \item Which disguises itself as legitimate software?
    \item Which commonly encrypts data?
    \item Which secretly collects information?
\end{enumerate}

\textbf{Answers:}

1. Virus.

2. Worm.

3. Trojan.

4. Ransomware.

5. Spyware.

% ============================================================
% TASK 1 SUMMARY
% ============================================================

\subsection{Task 1 Summary}

Viruses infect files, worms spread between systems, Trojans disguise malicious software, ransomware can restrict access to data, and spyware can secretly collect information. Understanding these types helps users recognize threats and helps cybersecurity professionals investigate malicious activity.

% ============================================================
% TASK 2
% ============================================================

\clearpage

\section{Task 2: Understanding Phishing}

\subsection{Introduction}

Phishing is a social-engineering technique in which attackers use deceptive messages, calls, or websites to trick people into revealing sensitive information, clicking harmful links, downloading malicious content, or performing an action that benefits the attacker.

% ============================================================
% WHAT IS PHISHING
% ============================================================

\subsection{What is Phishing?}

\subsubsection{Definition}

Phishing uses deceptive communication to trick a victim into providing sensitive information or performing an unsafe action. Attackers commonly impersonate trusted people, companies, services, or organizations.

\subsubsection{How a Phishing Attack Works}

A typical attack begins when an attacker prepares a deceptive message and delivers it through email, SMS, a phone call, a website, or another communication channel. The victim is encouraged to click, respond, enter credentials, make a payment, or provide sensitive information.

\begin{figure}[H]
\centering
\includegraphics[width=.72\textwidth]{images/phishing_attack_flow.png}
\caption{Phishing Attack -- General Process}
\end{figure}

\subsubsection{Why Do Attackers Use Phishing?}

Attackers use phishing to exploit human decision-making. Successful phishing may steal credentials, obtain personal or financial information, gain account access, deliver malware, redirect users to fraudulent sites, or convince victims to make payments.

\refs
{https://www.nist.gov/itl/smallbusinesscyber/guidance-topic/phishing}
{https://nvlpubs.nist.gov/nistpubs/TechnicalNotes/NIST.TN.2276.pdf}
{https://www.youtube.com/results?search_query=Professor+Messer+phishing}

% ============================================================
% EMAIL PHISHING
% ============================================================

\subsection{Email Phishing}

\subsubsection{Definition}

Email phishing uses deceptive emails to impersonate trusted organizations, companies, or individuals and persuade the recipient to perform an unsafe action.

\subsubsection{How It Works}

The attacker sends an email designed to appear legitimate. It may contain a malicious link, attachment, fake login page, payment request, or request for sensitive information.

\begin{figure}[H]
\centering
\includegraphics[width=.70\textwidth]{images/email_phishing.png}
\caption{Email Phishing -- Deceptive Email}
\end{figure}

\subsubsection{Real-World Example}

A user may receive an email appearing to come from a bank claiming that immediate account verification is required and providing a fraudulent login link.

\subsubsection{Cybersecurity Connection}

Security teams may use email filtering, URL analysis, threat intelligence, and user reporting to identify suspicious messages.

\subsubsection{Security Risk and Protection}

\textbf{Risk:} A user may disclose credentials or download malicious content.

\textbf{Protection:} Verify the sender; inspect links; avoid unexpected attachments; verify requests through official channels; report suspicious emails.

\refs
{https://consumer.ftc.gov/articles/how-recognize-avoid-phishing-scams}
{https://www.nist.gov/itl/smallbusinesscyber/guidance-topic/phishing}
{https://www.youtube.com/results?search_query=Professor+Messer+phishing}

% ============================================================
% SMISHING
% ============================================================

\subsection{Smishing}

\subsubsection{Definition}

Smishing is phishing conducted through SMS or other text messaging services.

\subsubsection{How It Works}

An attacker sends a text that may appear to come from a bank, delivery company, service provider, or other trusted organization, often containing a link or request for information.

\begin{figure}[H]
\centering
\includegraphics[width=.70\textwidth]{images/smishing.png}
\caption{Smishing -- SMS Phishing}
\end{figure}

\subsubsection{Real-World Example}

A user may receive a text claiming that a package cannot be delivered until a payment or address confirmation is completed.

\subsubsection{Cybersecurity Connection}

Smishing is important to mobile security because users may trust messages received on personal devices.

\subsubsection{Security Risk and Protection}

\textbf{Risk:} A victim may open a malicious link or submit credentials or payment information.

\textbf{Protection:} Do not trust unexpected messages; avoid suspicious links; verify through the official app/site; report suspicious messages; delete after reporting.

\refs
{https://consumer.ftc.gov/articles/how-recognize-avoid-phishing-scams}
{https://www.nist.gov/itl/smallbusinesscyber/guidance-topic/phishing}
{https://www.youtube.com/results?search_query=Professor+Messer+smishing+phishing}

% ============================================================
% VISHING
% ============================================================

\subsection{Vishing}

\subsubsection{Definition}

Vishing, or voice phishing, is phishing conducted through telephone or voice communication.

\subsubsection{How It Works}

The attacker creates a believable urgent situation, such as a supposed security problem or payment issue, and asks the victim for sensitive information.

\begin{figure}[H]
\centering
\includegraphics[width=.70\textwidth]{images/vishing.png}
\caption{Vishing -- Voice Phishing}
\end{figure}

\subsubsection{Real-World Example}

A caller may pretend to be from a bank's security department and ask for an OTP, password, or card information.

\subsubsection{Cybersecurity Connection}

Vishing demonstrates how social engineering can operate without malware or technical exploitation. Security awareness helps users recognize impersonation.

\subsubsection{Security Risk and Protection}

\textbf{Risk:} Victims may reveal credentials, financial information, or authentication codes.

\textbf{Protection:} Never disclose passwords or OTPs during unexpected calls; end suspicious calls; use official contact numbers; do not trust caller ID alone; report suspected fraud.

\refs
{https://consumer.ftc.gov/all-scams/phishing-scams}
{https://www.nist.gov/itl/smallbusinesscyber/guidance-topic/phishing}
{https://www.youtube.com/results?search_query=Professor+Messer+social+engineering}

% ============================================================
% FAKE WEBSITE / LINK PHISHING
% ============================================================

\subsection{Fake Website / Link Phishing}

\subsubsection{Definition}

Fake website or link phishing uses deceptive URLs or fraudulent websites that imitate legitimate services to collect credentials or other sensitive data.

\subsubsection{How It Works}

A deceptive link is distributed through email, SMS, social media, or another channel. The victim is directed to a fraudulent website that may visually imitate a legitimate service.

\begin{figure}[H]
\centering
\includegraphics[width=.72\textwidth]{images/fake_website_phishing.png}
\caption{Fake Website -- Link Phishing}
\end{figure}

\subsubsection{Real-World Example}

A phishing message may contain a link that appears to belong to an online service but leads to a fraudulent login page.

\subsubsection{Cybersecurity Connection}

Security teams can analyze URLs, domains, certificates, page content, and reported phishing indicators.

\subsubsection{Security Risk and Protection}

\textbf{Risk:} Entering credentials into a fraudulent website can give attackers access to an account.

\textbf{Protection:} Check the address carefully; avoid unexpected links; use known official websites/apps; enable multi-factor authentication; report fraudulent sites.

\refs
{https://www.nist.gov/itl/smallbusinesscyber/guidance-topic/phishing}
{https://nvlpubs.nist.gov/nistpubs/TechnicalNotes/NIST.TN.2276.pdf}
{https://www.youtube.com/results?search_query=Professor+Messer+phishing}

% ============================================================
% FIVE SIGNS OF PHISHING
% ============================================================

\subsection{Five Common Signs of a Phishing Message}

\begin{enumerate}
    \item \textbf{Suspicious sender:}
    The sender address may not match the claimed organization.

    \item \textbf{Urgent or threatening language:}
    The message pressures the recipient to act immediately.

    \item \textbf{Suspicious link:}
    The URL may not match the expected organization.

    \item \textbf{Request for sensitive information:}
    Passwords, financial information, or authentication codes may be requested.

    \item \textbf{Spelling, grammar, or unusual content:}
    Errors, unusual formatting, or unexpected requests can be warning signs.
\end{enumerate}

\begin{figure}[H]
\centering
\includegraphics[width=.98\textwidth]{images/phishing_warning_signs.png}
\caption{Five Common Signs of a Phishing Message}
\end{figure}

\refs
{https://www.nist.gov/itl/smallbusinesscyber/guidance-topic/phishing}
{https://nvlpubs.nist.gov/nistpubs/TechnicalNotes/NIST.TN.2276.pdf}
{https://www.youtube.com/results?search_query=NIST+phishing+awareness}

% ============================================================
% PHISHING COMPARISON
% ============================================================



\subsection{Phishing Types Comparison}

\begin{table}[H]
\centering
\small
\renewcommand{\arraystretch}{1.25}

\begin{tabularx}{\textwidth}{
|p{2.4cm}|p{2.2cm}|X|X|
}
\hline

\textbf{Type}
&
\textbf{Channel}
&
\textbf{Common Technique}
&
\textbf{Typical Target}
\\
\hline

Email Phishing
&
Email
&
Fake messages, links, attachments
&
Credentials and sensitive information
\\
\hline

Smishing
&
SMS / Text
&
Fake delivery, account, or payment messages
&
Mobile users and account information
\\
\hline

Vishing
&
Voice / Phone
&
Impersonation and urgent requests
&
Credentials, financial information, OTPs
\\
\hline

Fake Website / Link Phishing
&
Web / Links
&
Fraudulent websites and deceptive URLs
&
Login credentials and personal information
\\
\hline

\end{tabularx}

\caption{Comparison of Common Phishing Types}

\end{table}

% ============================================================
% PHISHING MINI PRACTICAL
% ============================================================

\subsection{Mini Practical / Try It Yourself}

A safe phishing-awareness activity can be performed without interacting with a real phishing website.

\begin{enumerate}
    \item Find a phishing-awareness example from NIST, FTC, or another recognized cybersecurity-awareness organization.
    \item Examine sender information.
    \item Identify suspicious links or requests without opening them.
    \item Identify urgency, threats, or unusual wording.
    \item Identify requests for sensitive information.
    \item Record the warning signs.
\end{enumerate}

\textbf{Safety Note:} Do not click suspicious links, download unknown attachments, submit personal information, or interact with real phishing websites.

% ============================================================
% PHISHING CYBERSECURITY CONNECTION
% ============================================================

\subsection{Cybersecurity Connection}

Phishing is important to cybersecurity because attackers can use social engineering to bypass technical controls by manipulating users. Security teams may use email security gateways, spam filters, URL analysis, threat intelligence, endpoint security tools, and user reporting.

% ============================================================
% PHISHING SECURITY RISKS
% ============================================================

\subsection{Security Risks and Protection}

\textbf{Risks:} Credential theft, financial fraud, account compromise, malware infection, identity theft, and unauthorized access.

\textbf{Protection:} Stop and verify unexpected requests; do not click suspicious links or attachments; use multi-factor authentication; verify through official channels; report suspicious messages; keep devices updated.

\refs
{https://consumer.ftc.gov/articles/how-recognize-avoid-phishing-scams}
{https://www.nist.gov/itl/smallbusinesscyber/guidance-topic/phishing}
{https://www.youtube.com/results?search_query=NIST+phishing+awareness}

% ============================================================
% SAFE RESPONSE
% ============================================================

\subsection{Safe Response to a Suspicious Message}

\begin{figure}[H]
\centering
\includegraphics[width=.62\textwidth]{images/phishing_protection.png}
\caption{Safe Response to a Suspicious Message}
\end{figure}

% ============================================================
% QUICK CHECK
% ============================================================

\subsection{Quick Check}

\begin{enumerate}
    \item What is phishing?
    \item Why do attackers use phishing?
    \item What is the difference between phishing and smishing?
    \item What does vishing use?
    \item Name five warning signs of phishing.
\end{enumerate}

\textbf{Answers:}

1. A social-engineering attack using deceptive communication.

2. To obtain information, access, money, or deliver malicious content.

3. Smishing specifically uses SMS/text messages.

4. Voice or telephone communication.

5. Suspicious sender, urgent language, suspicious link, request for sensitive information, and unusual spelling/grammar/content.

% ============================================================
% TASK 2 SUMMARY
% ============================================================

\subsection{Task 2 Summary}

Phishing is a social-engineering technique that manipulates users into clicking harmful links, opening malicious content, or revealing sensitive information. Email phishing, smishing, vishing, and fake website/link phishing use different channels but share the objective of deceiving the victim. Recognizing common warning signs helps users avoid unsafe actions.

% ============================================================
% END
% ============================================================

\end{document}